\section{Lexicon Construction and Representational Similarity}
\subsection{Lexical inventory and provenance}
The active dataset contains 613 lexical entries. Each entry identifies a word label and its corresponding concept, lemma, and phonological form nodes. The conceptual graph contains 1,205 nodes, including 593 nodes classified as deep-layer nodes in the source data; the phonological graph contains 1,616 nodes, including 1,003 deep-layer nodes. Deep-layer nodes supply representational structure rather than additional retrieval candidates; the underlying representational organization is described by \citet{stille2020}. One conceptual lexical root also belongs to the deep-layer category, so the conceptual counts are not disjoint categories that can simply be added.

The graphs were extracted from the existing WWT source data \citep{stille2020}. The current dataset documents 13 restored phonological mappings and one component correction for the entry \emph{Verbringen}. Original and restored links are recorded explicitly. These repairs are changes introduced in the current dataset relative to the source files used by Katharina Stille in the earlier modeling work \citep{stille2020}.

In this paper, eight words are used for detailed display and as baseline inputs: \emph{Tasse}, \emph{Tasche}, \emph{Kanne}, \emph{Katze}, \emph{Teller}, \emph{Topf}, \emph{Stuhl}, and \emph{Schrank}. This selection keeps the user interface and the detailed presentation manageable while providing a compact set of examples to illustrate how the model operates. The underlying model uses the full lexicon of 613 entries, all of which remain available as retrieval candidates in the associative memories.

\subsection{Recursive graph encoding}
\label{sec:graph-encoding}
Let $E(x)$ be the outgoing relation--target pairs of node $x$, $\mathbf b_x$ its random base vector, and $\mathbf r_\rho$ the vector assigned to relation $\rho$. A node is encoded as
\begin{align}
 \mathbf u_x &= \mathbf b_x+\sum_{(\rho,y)\in E(x)}\mathbf r_\rho\circledast\mathbf v_y,\\
 \mathbf v_x &= \frac{\mathbf u_x}{\norm{\mathbf u_x}}.
 \label{eq:recursive}
\end{align}
For indices $k=0,\ldots,D-1$, circular convolution is
\begin{equation}
 (\mathbf a\circledast\mathbf b)_k=\sum_{m=0}^{D-1}a_m b_{(k-m)\bmod D}.
\end{equation}
A leaf has no outgoing relations and is represented by its normalized base vector. Non-leaf nodes incorporate the recursively constructed target vectors, not separate target identity vectors. Each node is normalized after summation. The implementation checks for directed cycles because this recursive procedure requires an acyclic dependency structure.

Base and relation vectors are obtained by normalizing independent Gaussian draws. They are ordinary random unit vectors; no unitary-convolution constraint is imposed. A \emph{vector seed} fixes the pseudorandom choices from which the vectors are generated. In this implementation, each vector's random stream is derived deterministically from the seed, its representational level and role, and its node or relation identifier. The seed therefore specifies one numerical realization of the vector vocabulary, including the independently generated lemma vectors.

This choice makes the arbitrariness of the numerical encoding explicit. The lexical graph specifies which nodes and relations are present and which components are shared; the seed selects concrete vector values within that specification. For example, different seeds assign different numerical vectors to the handle-related component shared by \emph{Tasse} and \emph{Kanne}, while both words continue to share that component. The graph-defined relations are preserved across these realizations, and the exact similarities between word vectors can vary. The resulting geometry reflects both the specified lexical structure and its particular numerical realization.

Holding the graph, vector dimensionality, encoding procedure, identifiers, and software settings fixed, the same seed reproduces the same vector vocabulary. Fixing the seed thus preserves one chosen configuration for repeated inspection or comparisons between model conditions. Deliberately changing it generates alternative configurations of the same lexical specification; recording each alternative seed makes those configurations reproducible as well. This provides a way to examine how retrieval depends on the numerical encoding. The role of the seed is to control that choice; the linguistic adequacy of the resulting similarities remains a question for independent empirical comparison.

Because the random streams are tied to vector identifiers, existing vectors also remain stable when unrelated entries are added, provided that their dependency graphs remain unchanged. This supports comparisons between the eight-entry fixture and the 613-entry dataset with identical vectors for the shared entries. The previous WWT vectors used in the model developed by Katharina Stille were generated with a different randomization procedure.

\subsection{Conceptual similarity}
Concept nodes use the relation \texttt{AssocWith}. For example, the source graph associates \emph{Tasse} with both a dish-related node and a handle-related node, whereas \emph{Kanne} includes the handle-related node. Their encodings therefore contain a common bound contribution. The resulting overlap depends on all contributions and the normalization in Eq.~\eqref{eq:recursive}; it is not assigned directly as a fixed pairwise similarity value.

Shared components make similarity interpretable in terms of the conceptual relations encoded in the graph. Inspecting these relations helps explain why particular word vectors overlap and which aspects of their conceptual structure contribute to that overlap. For example, the source concept for \emph{Tasche} includes a connection to a vending-machine-related visual node, which contributes to its vector through the same encoding rule. The resulting similarities thus provide a way to examine how the associations specified in the lexical dataset shape the geometry of the conceptual representations.

\subsection{Phonological similarity}
Phonological form nodes use the relation \texttt{InclPhon} to incorporate recursively represented components. For example, \emph{Tasse} and \emph{Tasche} share the source component \texttt{P\_St\_ta}, while \emph{Kanne} and \emph{Katze} share \texttt{P\_St\_ka}. Common components at deeper levels can also contribute to overlap. Phonological similarity in this model thus expresses the overlap between vector representations constructed from the components and relations specified in the phonological graph. Shared structure contributes to this overlap through the recursive encoding procedure, with the resulting similarity determined by the complete composition and normalization of the word vectors. The measure therefore provides an explicit link between the phonological structure represented in the lexicon and the similarity values used to interpret lexical retrieval.

To construct a phonological representation, the model first binds each component vector to the vector for its relation, such as \texttt{InclPhon}, and then adds these contributions to the node's base vector. For example, if a node includes two components with bound contributions $\mathbf q_1$ and $\mathbf q_2$, their combined contribution is $\mathbf q_1+\mathbf q_2$. Addition gives the same result when these contributions are processed in the reverse order. The resulting representation therefore depends on the components and relations specified in the graph and remains stable across different orderings of the same list of graph connections. Position within a word---for example, whether a component belongs to the first or second syllable---is a separate aspect of phonological structure. Encoding such a distinction requires it to be expressed in the component identities or relations supplied to the encoding procedure. The graph thus specifies the phonological content and structural distinctions that are combined into a static vector representation for lexical retrieval. These structural distinctions describe positions within a word; they do not specify when its components become active or how long they remain active during articulation. The present encoding therefore does not generate a temporal sequence of syllables or articulatory gestures.

\subsection{Within-level similarity and cross-level association}
Within-level similarity describes the overlap between lexical representations on the same representational level. It applies to all three levels: conceptual, lemma, and phonological. In the following notation, $L$ denotes any one of these levels. For normalized vectors representing lexical entries $i$ and $j$ at the selected level, define
\begin{equation}
 S^{(L)}_{ij}=(\mathbf v_i^{(L)})^{\mathsf T}\mathbf v_j^{(L)}.
 \label{eq:similarity}
\end{equation}
This quantity is both a dot product and a cosine similarity. Applied separately to the conceptual vectors $\mathbf c_i$, lemma vectors $\boldsymbol\ell_i$, and phonological form vectors $\mathbf p_i$, it yields three similarity matrices, each describing the static representational geometry of its level. Values near zero indicate small overlap. Negative values express opposing vector components in this geometric measure; their interpretation in terms of linguistic relations or neural activity requires additional evidence. Independent random unit vectors have expected dot product zero and variance $1/D$, providing a reference for incidental overlap. The distribution of similarities among structured vectors additionally reflects their encoded relations.

At the conceptual and phonological levels, the graph-based construction introduces shared structure into the vectors. Conceptual overlap can reflect common associations, while phonological overlap can reflect common components of word forms. The lemma level currently represents lexical identity using independently generated vectors. Similarity is measured there by the same equation, with the resulting overlap arising from the random-vector construction. These lemma vectors are approximately orthogonal, providing distinct identities for the mappings between concepts and phonological forms.

A possible extension would enrich lemma representations with shared grammatical features. For example, lemmas belonging to the same word class, such as nouns or verbs, could include a common contribution encoding that class alongside their individual identity contributions. Other grammatical properties could be incorporated in the same way. Such a construction would introduce systematic lemma-level similarities based on grammatical structure. This extension would require an explicit specification of the features and their encoding; the present implementation uses the independent identity vectors described above.

Cross-level association links the representations of the same lexical entry across levels. The model specifies both concept--lemma associations through the pairs $(\mathbf c_i,\boldsymbol\ell_i)$ and lemma--phonological form associations through the pairs $(\boldsymbol\ell_i,\mathbf p_i)$. In production, these associations support retrieval from concept to lemma and from lemma to phonological form; in perception, they support retrieval in the reverse direction. Their correspondence is established by the shared lexical entry index $i$. Thus, within-level similarity describes the geometry among representations on one level, while cross-level association specifies which representation on another level belongs to the same lexical entry. The associative memories implement these correspondences between the independently constructed vector spaces.
